% BHU PYQ -- Manually transcribed & error-corrected
% Paper: BSC-04A: Elements of Geology - II
% Exam: B.Sc. (Hons.) Semester IV Examination, 2023-24

\documentclass[11pt,a4paper]{article}
\usepackage[a4paper,top=2.5cm,bottom=2.5cm,left=2.8cm,right=2.8cm,headheight=14pt]{geometry}
\usepackage{amsmath,amssymb}
\usepackage[T1]{fontenc}
\usepackage[utf8]{inputenc}
\usepackage{lmodern,microtype,fancyhdr,enumitem,hyperref,lastpage,parskip}

\hypersetup{colorlinks=true,urlcolor=black,linkcolor=black,
  pdftitle={BSC-04A Elements of Geology - II -- BHU B.Sc. Sem IV 2023-24}}
\pagestyle{fancy}\fancyhf{}
\renewcommand{\headrulewidth}{0.4pt}\renewcommand{\footrulewidth}{0.4pt}
\fancyhead[L]{\small   \textit{Banaras Hindu University}}
\fancyhead[R]{\small   \textit{BSC-04A: Elements of Geology - II}}
\fancyfoot[C]{\small Page  \thepage\ of \pageref{LastPage}}


\newlist{parts}{enumerate}{2}
\setlist[parts,1]{label=(\alph*),leftmargin=*,itemsep=5pt,topsep=3pt}

\newlist{romanparts}{enumerate}{2}
\setlist[romanparts,1]{label=(\roman*),leftmargin=*,itemsep=5pt,topsep=3pt}

\newcommand{\pts}[1]{\hfill   \textit{[#1]}}


\begin{document}

\begin{center}
  {\large\bfseries BANARAS HINDU UNIVERSITY}\\[2pt]
  {\normalsize B.Sc. (Hons.) --- Semester IV Examination, 2023-24}\\[6pt]
  \rule{0.8\linewidth}{0.5pt}\\[6pt]
  {\large\bfseries Paper: BSC-04A --- Elements of Geology - II}\\[4pt]
  {\normalsize Time: 3 Hours \hfill Maximum Marks: 70}\\[2pt]
  {\small Ref. No.: 074447}
\end{center}
\rule{\linewidth}{0.4pt}
\smallskip



\noindent   \textit{Instructions: Answer   \textbf{five} questions in all, selecting at least   \textbf{two} questions from each of Sections A and B. Section-C is compulsory. All questions carry equal marks.}
\bigskip

\bigskip


\noindent {\large\bfseries SECTION --- A}
\rule{\linewidth}{0.4pt}\smallskip




\noindent   \textbf{Question 1.} Draw a geological time scale marking the boundaries of EON and ERA in million years. Explain the evolution of life through different geological periods. \pts{14}

\medskip


\noindent   \textbf{Question 2.} Answer the following:

\begin{parts}
  \item How is the physiography of the Northern Mountains different from that of the Peninsular Plateaus? \pts{7}
  \item What are Steno's Laws and how do they help in establishing the relative dating of different stratigraphic units? \pts{7}
\end{parts}

\medskip


\noindent   \textbf{Question 3.} Answer the following:

\begin{parts}
  \item Discuss the differences between lithostratigraphic and biostratigraphic units. \pts{5}
  \item Discuss the desertification of Thar and mention the characteristic topography. \pts{5}
  \item Comment on the `Great Dying Event'. \pts{4}
\end{parts}

\medskip


\noindent   \textbf{Question 4.} Answer the following:

\begin{parts}
  \item Name the eastern and western coastal plains of India. \pts{4}
  \item Explain the Law of Uniformitarianism. \pts{4}
  \item Describe the division of the Northern Plains. \pts{3}
  \item Explain the physiographic divisions of India. \pts{3}
\end{parts}

\bigskip


\noindent {\large\bfseries SECTION --- B}
\rule{\linewidth}{0.4pt}\smallskip




\noindent   \textbf{Question 5.} Discuss the Siwalik vertebrate fauna and describe their significance in paleoclimatic interpretations. \pts{14}

\medskip


\noindent   \textbf{Question 6.} Define a fossil. Describe different modes of occurrence of fossils. \pts{14}

\medskip


\noindent   \textbf{Question 7.} Write explanatory notes on the following:

\begin{parts}
  \item Conditions for fossilization. \pts{6}
  \item Scope of Palaeontology. \pts{4}
  \item Index fossil and its use. \pts{4}
\end{parts}

\bigskip


\noindent {\large\bfseries SECTION --- C}
\rule{\linewidth}{0.4pt}\smallskip




\noindent   \textbf{Question 8.} Define the following in a maximum of five sentences each: \pts{$2  \times 7 = 14$}

\begin{romanparts}
  \item Unconformity
  \item Adam's Bridge
  \item Minicoy Island
  \item Kankar
  \item Living fossil
  \item Fossil in toto
  \item Silicification
\end{romanparts}

\vfill\rule{\linewidth}{0.4pt}

\begin{center}\small
\href{https://bhu-exam.vercel.app/}{Click here to visit our website}\\[4pt]
\href{https://me-aryan.vercel.app/}{Click here to visit our contributor}\\[4pt]
\href{https://quashan-me.vercel.app/}{Click here to visit our contributor}
\end{center}
\end{document}
